% Fuentes: controles variacionales EH y cosmológico, tras sus deducciones.
\begin{frame}[label=nav-frame-040]{Casos de referencia: ecuaciones y frontera}
\TablaSetup
\begin{tabular}{@{}p{.24\textwidth}p{.44\textwidth}p{.25\textwidth}@{}}
\TablaCabecera{Objeto} & \TablaCabecera{Einstein--Hilbert} & \TablaCabecera{Cosmológico} \\
\TablaLinea
Derivada del momento & $\nabla_mP_{\mathrm{EH}}^{abcd}=0$ & $P_\Lambda^{abcd}=0$ \\[2mm]
Tensor de campo & $E_{ab}^{(\mathrm{EH})}=G_{ab}$ & $E_{ab}^{(\Lambda)}=\Lambda g_{ab}$ \\[2mm]
Identidad de Bianchi & $\nabla^aG_{ab}=0$ & $\nabla^aE_{ab}^{(\Lambda)}=0$ \\[2mm]
Frontera & $\begin{aligned}[t]\delta v_{\mathrm{EH}}^j={}&g_{ab}\nabla^j\delta g^{ab}\\[-1mm]&-\nabla_a\delta g^{aj}\end{aligned}$ & $\delta v_\Lambda^j=0$ \\
\TablaLinea
\end{tabular}\par
\vspace{4mm}

El término cosmológico contribuye al volumen mediante $\Lambda g_{ab}$,
aunque sus momentos y su frontera sean nulos.
\end{frame}
